% Transcription — fonds Alexandre Grothendieck, shelfmark 40, batch 1 (pages 1-20).
% Established from the facsimile archives/batches/40/batch-01.pdf, cut from a
% copy of the folder fetched through the project's own relay
% (https://grothendieck-relay.onrender.com/source/40.pdf), not uploaded by the
% user; PDF page k is the archivists' page k, with no cover sheet. Per the
% skill transcribe-grothendieck. The folder is twenty-seven pages; this is the
% first of its two batches (pages 21-27 are batch 2, transcribed in parallel
% by another pass and already in transcripts/40/batch-02.fr.tex, whose
% notation this file follows).
%
% Pages skipped: none. Pages summarised: none; the diagrams of page 13 are
% given in part (first and last column of each computation), the middle
% columns being described in a \note. All twenty pages are his hand, blue
% ink, fast drafting; page 15 carries only three lines. No pagination of his
% is visible and no date. Page 1 carries, top left, a diagonal title in red
% crayon, « Calculs sur des fibrés principaux affines » (the last word under
% red strokes), and a circled « Hopf » in black ink; whether the red crayon is
% his hand cannot be told from the scan. It is very probably what the
% archivists' title « Calculs [affines sur des] fibrés principaux affines »
% was taken from; their bracketed words are theirs, not his, and the heading
% is given here as a \marginal with a \note, not as a title.
%
% What the batch is about. Affine group schemes over a commutative ring A
% and their principal homogeneous spaces (« fibrés principaux homogènes »),
% worked out on algebras. P. 1: a group in the dual of the category of
% commutative A-algebras — B with diagonal p_B : B -> B (x)_A B, augmentation
% epsilon (J = Ker epsilon), antipode x -> x-check, with its properties;
% when B is free of finite rank, the dual structure on U. P. 2: U as a
% unital coalgebra with Delta_U, eta_U, an associative unital algebra
% structure compatible with it, and the antipode; then an A-algebra C with
% B as « groupe d'opérateurs à gauche ». Pp. 3-4: the coaction q_C : C ->
% B (x) C (associativity and unit diagrams, q_C(y) in 1 (x) y + J (x) C);
% with B free of finite rank, q_C as q_U : U -> Hom_A(C, C'), and C as a
% U-module algebra (u(xy) = sum (v_i x)(w_i y)). Pp. 4-5 (much cancelled):
% C principal homogeneous iff C (x) C -> B (x) C, (q_C, p_2), is an
% isomorphism, iff there is g_C : B -> C (x) C making (*) and (**) commute.
% Pp. 5-8: the quotient of C by B in the category is the subalgebra C' of
% invariants (« primitifs », then « invariants »); a warning (p. 7) that the
% quotient in the affine category can differ from that in schemes (E/G for
% G a Borel subgroup: a point); for C principal homogeneous and A a direct
% factor of C, C' = A; the N.B. of p. 8, C = B/JB, where the invariants are
% A/J. P. 9: with C free, phi : C (x) C -> Hom(U, C) is an isomorphism iff
% psi : C (x) U -> Hom(C, C), psi(x (x) u).y = x u(y), is one (margin « ? »),
% with a cancelled base-change argument (A' = C). P. 10: the same in terms of
% U, with the convolution v * w on Hom_A(U, C). Pp. 11-15: the categorical
% definition (E formally principal homogeneous under G iff (q_E, p_2) is an
% isomorphism, iff g_E exists with (i), (ii)); translated into algebras;
% psi : Hom_A(B, C) -> Hom_A(C, C), psi(v) = Delta_C (v (x) id_C) q_C, is an
% isomorphism if phi is, with inverse psi'(w) = Delta_C (w (x) id_C) g_C,
% verified by the column diagrams of p. 13; the converse (pp. 14-15) when C
% has a finite basis, via rho : Hom(Hom(C, C), C) -> C (x) C, left off at
% « Pour ceci, ---- » with « ??? » in the margin. Pp. 16-20: the example
% that batch 2 continues — the hyperalgebra U = A[X]/X^p A[X] in
% characteristic p (Delta X = X (x) 1 + 1 (x) X, epsilon(X) = 0, X-check =
% -X), dual of the group B of batch 2; an action of U on C is an A-derivation
% D with D^p = 0; C principal homogeneous iff the D^i (0 <= i <= p-1) form a
% C-basis of Hom_A(C, C); for C = A[u]/(u^p - a), Du = 1, the matrix of the
% D^i, triangular with determinant 0!1!...(p-1)!, gives the « Proposition »
% (p. 18): an A-algebra with a one-element p-basis is principal homogeneous
% under U; « Questions » (hyperalgebras in Hom_A(C, C), restricted
% enveloping algebras); p. 19, « 2) » the classification of principal
% homogeneous spaces under B begins: the elements killed by D are A, via a
% lemma on direct-factor submodules M and their orthogonals M° in
% Hom(C, C); p. 20: existence and uniqueness mod A of u with Du = 1, the
% local argument (image of D free of rank p-1, containing A), x^p in A, and
% « si Du = 1 les u^i forment une base de C sur A » — continued on p. 21 of
% batch 2 by the struck proof « A est local. Montrons d'abord que les u^i
% ... sont libres sur A ».
%
% Notation fixed once for the batch, aligned with batch 2. The group algebra
% is B (plain, as on pp. 21-23 of batch 2); the category is \mathcal{C}. His
% U with a curled stroke, the dual hyperalgebra, is set \mathcal{U}, as batch
% 2 sets \mathcal{U}' on p. 23. The script letter he uses on pp. 5-6 for a
% second algebra (« h : D -> C ») is \mathcal{D}. His Delta_C on pp. 10-13
% is the multiplication C (x) C -> C and is kept so; Delta_U is the
% diagonal of U. eta_U is his augmentation of U (a hooked eta), epsilon that
% of B. The antipode is \check{x}. Tensor products are over A unless he
% writes otherwise. On p. 4 he writes q_U in the first condition and p_U in
% the second: kept, with a note. His underlinings are \emph.
%
% Readings to check first: the red heading on p. 1 and whether « affines »
% is struck; the first words of p. 2 (« D\ill{} de cogèbre ») and of p. 3;
% the whole cancelled block at the top of p. 5 and the closing sentence of
% p. 5 (« Nous \ill{} ... si C soit fibré principal ... »); « on voit qu'on a
% tout ce qu'il faut » on p. 6; the first paragraph of p. 7 (the example
% with the Borel subgroup: « groupe linéaire », « réduit au point »); the
% bracket of p. 8 (« ayant une base contenant 1_C »); the cancelled lower
% half of p. 9; the word before « e » on p. 11 (left \ill{}); « on démontre » on p. 12; the
% « Questions » of p. 18; the margins and the lemma of p. 19; the framed
% local argument and the rotated left margin of p. 20.
% Apparatus: 109 \ill{}, 45 \uncertain{}.
%
% Pass: Opus 5.5 (claude-opus-5-5), 2026-09-26 — first pass, unchecked against the pages by a human.

\documentclass[11pt,a4paper]{article}
% Le préambule commun à toutes les transcriptions du fonds Grothendieck.
%
% Il tient en une idée : **l'incertitude se compose, elle ne se résout pas.**
% Une transcription qui lisse un mot douteux a détruit la seule chose pour
% laquelle on la faisait — un lecteur doit pouvoir savoir, à chaque endroit,
% ce qui était sur le feuillet et ce qui a été ajouté. Les six macros qui
% suivent sont donc l'appareil critique, et non de la décoration.
%
% Les mêmes commandes sont reconnues par `scripts/render.mjs`, qui produit la
% vue de lecture du site. Une macro ajoutée ici doit l'être aussi là-bas,
% faute de quoi le rendu échouera bruyamment — ce qui est voulu : un rendu
% silencieusement faux est pire qu'un rendu absent.

\ProvidesPackage{grothendieck}[2026/08/08 Transcriptions du fonds Grothendieck]

\usepackage{amsmath,amssymb,amsthm}
\usepackage{mathrsfs}
\usepackage{stmaryrd}
% Les diagrammes commutatifs sont partout dans ce fonds : c'est la langue
% même de la géométrie algébrique de Grothendieck, et une transcription qui
% les rendrait en prose ne transcrirait rien.
\usepackage{tikz-cd}
% Français par défaut — c'est la langue des feuillets — mais l'anglais est
% chargé aussi : la traduction et le résumé sont des documents anglais, et la
% typographie française (espaces avant les deux-points) y serait une faute.
% Ces documents basculent par \selectlanguage{english} en tête de corps.
\usepackage[english,french]{babel}
\usepackage{fontspec}
\usepackage{geometry}
\geometry{a4paper,margin=25mm,marginparwidth=28mm}
\usepackage{marginnote}
\usepackage{xcolor}
% ulem, not soul. `soul`'s \st and \ul are built on a character-by-character
% scan that XeLaTeX's font machinery breaks: the document fails with an
% undefined control sequence rather than degrading. ulem does the same two jobs
% and is Unicode-engine safe. `normalem` stops it hijacking \emph, which the
% transcriptions use for Grothendieck's own emphasis.
\usepackage[normalem]{ulem}
\usepackage{hyperref}
\hypersetup{colorlinks=true,linkcolor=black,urlcolor=black,pdfborder={0 0 0}}

% --- Métadonnées du lot -----------------------------------------------------
% Reprises telles quelles par le script de rendu, qui les affiche en tête de
% la vue de lecture. Elles fixent ce que le fichier prétend couvrir : sans
% elles, une transcription flotte sans point d'attache dans un fonds de seize
% mille feuillets.

\newcommand{\folder}[1]{\gdef\grothFolder{#1}}
\newcommand{\batch}[1]{\gdef\grothBatch{#1}}
\newcommand{\leaves}[2]{\gdef\grothFirst{#1}\gdef\grothLast{#2}}
\newcommand{\foldertitle}[1]{\gdef\grothFoldertitle{#1}}
\gdef\grothFolder{?}\gdef\grothBatch{?}\gdef\grothFirst{?}\gdef\grothLast{?}\gdef\grothFoldertitle{}

% --- Le résumé --------------------------------------------------------------
%
% Il ouvre la lecture modernisée, sous le titre « Résumé », et il est composé
% en petit. Ce n'est pas un ornement : c'est le seul passage du document qui
% ne soit pas des mathématiques, et un lecteur doit pouvoir le reconnaître
% avant de l'avoir lu — puis le sauter s'il connaît déjà le sujet, ou s'y
% arrêter s'il ne le connaît pas. Le filet qui le suit marque la reprise du
% corps.

\newenvironment{resume}
  {\small\setlength{\parskip}{0.45em}}
  {\par\medskip\hrule\medskip}

% --- Mots-clés ---------------------------------------------------------------
%
% En anglais, en clôture du résumé de la lecture modernisée : le vocabulaire
% moderne sous lequel chercher ce que les feuillets construisent. Le manifeste
% du site (`npm run manifest`) les extrait pour étiqueter la cote entière —
% cette ligne est la source unique des tags, il n'y a pas de fichier à côté.
\newcommand{\keywords}[1]{\par\smallskip\noindent
  {\footnotesize\textsc{Keywords} --- \textit{#1}}\par}

% --- L'appareil critique ----------------------------------------------------

% Le numéro de feuillet, en marge. C'est l'ancre qui permet de régler un
% désaccord sur une formule en regardant *un* feuillet plutôt qu'une chemise
% de six cents. Le site s'en sert aussi pour tourner les pages du fac-similé
% quand on fait défiler la transcription.
\newcommand{\leaf}[1]{%
  \marginnote{\footnotesize\color{gray!70}\textsf{#1}}%
  \label{leaf:#1}\ignorespaces}

% Illisible : on ne devine pas. Un mot inventé qui se lit comme les autres est
% le pire résultat possible.
\newcommand{\ill}{\textcolor{red!60!black}{[\,\ldots\,]}}

% Lecture douteuse : le mot est donné, et signalé comme incertain.
\newcommand{\uncertain}[1]{\uline{#1}}

% Ajout de l'éditeur — une lettre manquante, un mot sous-entendu.
\newcommand{\add}[1]{\textcolor{blue!50!black}{[#1]}}

% Biffé par Grothendieck lui-même. Ce qu'il a rayé fait partie du document :
% une rature dit souvent où la pensée a changé de direction.
\newcommand{\struck}[1]{\sout{#1}}

% Note du transcripteur, sur l'état matériel du feuillet.
\newcommand{\note}[1]{\par\smallskip{\small\color{gray!60!black}\textit{[#1]}}\par\smallskip}

% Note marginale de Grothendieck, distincte du corps du texte.
\newcommand{\marginal}[1]{\par\smallskip\noindent\hspace{1em}%
  {\small\color{gray!60!black}#1}\par\smallskip}

% --- Titre ------------------------------------------------------------------

\AtBeginDocument{%
  \begin{center}
    {\large\bfseries Fonds Alexandre Grothendieck --- cote n\textsuperscript{o}~\grothFolder}\\[2pt]
    {\itshape\grothFoldertitle}\\[4pt]
    {\small feuillets \grothFirst--\grothLast\ (lot \grothBatch)}\\[2pt]
    {\footnotesize\color{gray!60!black}Transcription établie d'après le fac-similé numérisé
     par l'Université de Montpellier.\\
     Les lectures incertaines sont soulignées, les illisibles notées [\,\ldots\,],
     les ajouts entre crochets.}
  \end{center}
  \vspace{1em}}

\endinput


\folder{40}
\batch{1}
\pages{1}{20}
\dating{[années 1960-1970]}
\watermark{Édition de démonstration}
\foldertitle{Calculs [affines sur des] fibrés principaux affines : notes manuscrites (s.d.)}

\begin{document}

\page{1}

\marginal{Calculs sur des fibrés principaux \uncertain{affines}}\note{au
crayon rouge, en diagonale dans le coin supérieur gauche, le dernier mot
sous des traits rouges ; on ne peut dire d'après la reproduction si ce
crayon est de sa main. Tout près, à l'encre noire, « Hopf » cerclé, et deux
traits obliques qui séparent le coin du texte}

On se place dans la catégorie duale de la catégorie des algèbres $B$
(commutatives avec unité) sur l'anneau comm. $A$. Une telle $B$ est un
groupe de la catégorie si c'est une \emph{$A$-algèbre commutative unitaire},
s'est donné un \emph{hom. d'algèbres}
\[
p_B : B \longrightarrow B \otimes_A B
\]
(« application diagonale ») qui soit « \emph{associative} »,\note{« d'algèbre »
est écrit sous « associative », souligné avec lui} pour laquelle il
\emph{existe} une \emph{augmentation} $\varepsilon : B \to A$
\struck{\ill{}} telle que
\[
\delta_B x = p_B x - (x \otimes 1 + 1 \otimes x) \in J \otimes J
\]
(où $J = \operatorname{Ker} \varepsilon$) --- cette augmentation est
alors unique --- et un \emph{hom d'algèbres} $x \mapsto \check{x}$
satisfaisant à la condition
\[
x \in B, \quad p_B x = \sum y_i z_i \;\Longrightarrow\;
\sum \check{y}_i z_i = \sum y_i \check{z}_i = \varepsilon(x) . 1
\]
\uncertain{un} tel hom d'algèbres est alors unique, et est une involution.
De plus on a
\[
\varepsilon(\check{x}) = \varepsilon(x)
\]
si $p_B x = \sum y_i \otimes z_i$, alors
$p_B \check{x} = \sum_i \check{z}_i \otimes \check{y}_i$.

\emph{Si} $B$ est libre de rang fini sur \uncertain{$B$},\note{sic ; on
attend $A$} alors ces structures se \uncertain{reflètent} directement dans
$\mathcal{U}$ par la donnée \struck{d'une} \struck{loi} \struck{d'algèbre} \struck{(non} \struck{nécessairement} \struck{commutative)} \struck{associative} \struck{unitaire} \struck{sur} \struck{$A$,} \struck{telle} \struck{que} \struck{\ill{}} \struck{il} \struck{existe} \struck{\ill{}} \struck{$x \mapsto \check{x}$}\note{les deux dernières lignes de la page sont
encadrées et barrées ; au-dessus de la dernière, un ajout biffé lui aussi,
\ill{}}

\page{2}

D\ill{} de cogèbre \add{unitaire} sur $A$, \struck{\ill{}} donnée
par\note{« unitaire », souligné, est ajouté au-dessus de la ligne avec un
renvoi}
\[
\Delta_{\mathcal{U}} : \mathcal{U} \longrightarrow \mathcal{U} \otimes_A \mathcal{U},
\qquad
\eta_{\mathcal{U}} : \mathcal{U} \longrightarrow A
\]
augmentation (surjective)
\struck{compatibles} \struck{avec} \struck{\ill{}} \struck{\ill{}} \struck{\ill{}} \struck{\ill{}} \struck{d'une} \struck{augmentation} \struck{$\delta_{\mathcal{U}} x$} \struck{${}= \Delta_{\mathcal{U}} x -{}$} \struck{\ill{}}\note{trois lignes
encadrées et barrées d'une ligne ondulée}
qui soit \emph{commutative}, \emph{associative}, d'une \emph{structure
d'$A$-algèbre} sur $\mathcal{U}$ qui soit \emph{associative},
\emph{unitaire} (non nécessairement commutative) et telle que
$\mathcal{U} \otimes_A \mathcal{U} \to \mathcal{U}$ soit un hom de cogèbres
unitaires [ou, ce qui revient au même, que $\Delta_{\mathcal{U}}$ est un hom
d'algèbres associatives unitaires]. Enfin on suppose qu'il existe un hom
\add{de cogèbres unitaires} $x \mapsto \check{x}$, ayant la propriété
suivante :\note{« de cogèbres unitaires », souligné, est ajouté au-dessus de
la ligne}
\[
x \in \mathcal{U}, \quad \Delta_{\mathcal{U}} x = \sum y_i \otimes z_i
\;\Longrightarrow\;
\sum y_i \check{z}_i = \sum \check{y}_i z_i = \eta_{\mathcal{U}}(x) . 1_{\mathcal{U}} .
\]
Alors $x \mapsto \check{x}$ est un antiautomorphisme involutif de la
structure d'algèbre de $\mathcal{U}$, et \struck{\ill{}} il est défini de
façon unique par les conditions ci-dessus.

Soit maintenant $C$ une $A$-algèbre. On veut \uncertain{expliciter} ce que
signifie que $C$ admet $B$ comme « groupe d'opérateurs à \struck{droite}
gauche » dans la catégorie $\mathcal{C}$. Cela signifie ceci :

\page{3}

\uncertain{On se donne} un \emph{hom. d'algèbres unitaires}
\[
q_C : C \longrightarrow B \otimes_A C
\]
satisfaisant aux conditions suivantes :

\emph{Associativité} (on a commutativité du diagramme

\begin{tikzcd}
B \otimes B \otimes C & B \otimes C \arrow[l, "\mathrm{id}_B \otimes q_C"'] \\
B \otimes C \arrow[u, "p_B \otimes \mathrm{id}_C"] & C \arrow[l, "q_C"] \arrow[u, "q_C"']
\end{tikzcd}

\emph{Unité} (on a commutativité du

\begin{tikzcd}
A \otimes C = C & B \otimes C \arrow[l, "\varepsilon \otimes \mathrm{id}_C"'] \\
 & C \arrow[ul, "\sim"] \arrow[u, "q_C"']
\end{tikzcd}

i.e. on a $q_C(y) \in 1 \otimes y + J \otimes C$ (où
$J = \operatorname{Ker} \varepsilon$).

Supposons encore que $B$ soit libre de rang fini, de sorte qu'on peut lui
substituer $\mathcal{U}$. Alors la donnée d'un $A$-hom
$q'_C : C \to B \otimes_A C'$ équivaut (pour tous les $A$-modules $C$, $C'$)
à un $A$-hom $q_{\mathcal{U}} : \mathcal{U} \to \mathrm{Hom}_A(C, C')$ [car
$B \otimes_A C' \simeq \mathrm{Hom}_A(\mathcal{U}, C')$].\note{sur l'indice
de $q_{\mathcal{U}}$, des lettres surchargées} Si $C$, $C'$ sont des
$A$-algèbres unitaires, alors dire que $q_C$ est un hom de $A$-algèbres
unitaires signifie que l'on a les conditions

\page{4}

\[
\left\lbrace
\begin{array}{l}
q_{\mathcal{U}}(y) . 1_C = \eta_{\mathcal{U}}(y) . 1_{C'} \\[4pt]
x, y \in C,\ u \in \mathcal{U},\ \Delta_{\mathcal{U}} u = \sum v_i \otimes w_i
\;\Longrightarrow\;
p_{\mathcal{U}}(u)(xy) = \sum \bigl(p_{\mathcal{U}}(v_i) x\bigr)\bigl(p_{\mathcal{U}}(w_i) y\bigr)
\end{array}
\right.
\]
\note{$q_{\mathcal{U}}$ dans la première ligne, $p_{\mathcal{U}}$ dans la
seconde, pour le même homomorphisme : tel sur la page}

Si enfin $C = C'$, alors \struck{\ill{}} les \uncertain{conditions}
d'associativité et d'unité signifient que \struck{$\mathcal{U}$ \ill{}}
$\mathcal{U} \to \mathrm{Hom}_A(C, C)$ est un \emph{hom de $A$-algèbres
unitaires}. Considérant alors $C$ comme un \emph{$\mathcal{U}$-module}
\add{unitaire} (i.e. \ill{}), les deux conditions \uncertain{précédentes}
s'écrivent ainsi :
\[
\left\lbrace
\begin{array}{l}
u . 1_C = \eta(u) . 1_C \\[4pt]
x, y \in C,\ u \in \mathcal{U},\ \Delta_{\mathcal{U}} u = \sum v_i \otimes w_i
\;\Longrightarrow\; u(xy) = \sum (v_i x)(w_i y)
\end{array}
\right.
\]

Nous voulons maintenant exprimer que $C$, considérée comme ayant un groupe
d'opérateurs $B$ dans la catégorie $\mathcal{C}$ des $A$-algèbres comm. avec
unité, est \struck{principal \ill{} \ill{}} \add{fibré} principal
\add{homogène} \struck{de base}\note{« fibré » et « homogène » sont écrits
au-dessus de la ligne, sur des mots biffés}
\struck{(i.e.} \struck{\ill{}} \struck{est} \struck{une} \struck{$A$-algèbre} \struck{commutative).} \struck{Cela} \struck{signifie} \struck{donc} \struck{qu'il} \struck{existe} \struck{un} \struck{hom} \struck{de} \struck{$A$-algèbres} \struck{unitaires} \struck{$C \otimes_A C \leftarrow
B$}\note{ces lignes sont encadrées et barrées de deux lignes ondulées ; le
cadre descend jusqu'à la flèche $C \otimes_A C \leftarrow B$}

\page{5}

\struck{On} \struck{suppose} \struck{donc} \struck{\ill{}} \struck{\ill{}} \struck{\ill{}} \struck{$A$-algèbres} \struck{\ill{}} \struck{que} \struck{$f : \mathcal{D} \to C$} \struck{\ill{}} \struck{un} \struck{hom} \struck{de} \struck{$A$-algèbres} \struck{unitaires} \struck{$g_C : B \to C \otimes C$} \struck{\ill{}} \struck{\ill{}} \struck{le} \struck{diagramme} \struck{suivant} \struck{$\leftarrow B \otimes C \xleftarrow{q_C} C$} \struck{ayant} \struck{la} \struck{propriété} \struck{suivante} \struck{:}\note{le haut de la page est barré de lignes horizontales et de
cinq longues diagonales ; lecture très incomplète}
l'hom composé
\[
B \otimes C \xleftarrow{\ \mathrm{id}_B \otimes \Delta_C\ } B \otimes C \otimes C
\xleftarrow{\ q_C \otimes \mathrm{id}_C\ } C \otimes C
\]
[i.e. l'hom $C \otimes C \xrightarrow{(q_C, p_2)} B \otimes C$] est un
\emph{isomorphisme} (de $A$-algèbres). Il revient au même de dire qu'il
\emph{existe} un \emph{hom de $A$-algèbres unitaires}
$g_C : B \to C \otimes C$ tel que l'on ait commutativité des

(*)

\begin{tikzcd}
C \otimes C & B \otimes C \arrow[l, "{(g_C, p_2)}"'] & C \arrow[l, "q_C"'] \arrow[ll, bend left=30, "p_1"]
\end{tikzcd}

(**)

\begin{tikzcd}
B \otimes C & C \otimes C \arrow[l, "{(q_C, p_2)}"'] & B \arrow[l, "g_C"'] \arrow[ll, bend left=30, "p_1"]
\end{tikzcd}

\note{dans (*), entre $C \otimes C$ et $B \otimes C$, un groupe biffé
illisible ; au-dessus, $(g_C, p_2)$ récrit sur une première étiquette ; une
boucle de crayon relie (**) au début de la ligne suivante}

Nous \ill{} \uncertain{examinons} \ill{} \ill{} les conditions, il est
\uncertain{bien} \uncertain{certain} que $C$ soit fibré principal
\uncertain{hom.} sous $G$, \ill{} \struck{\ill{}} \ill{} l'objet unité $A$
de $\mathcal{C}$. Supposons donc \ill{} un hom de $A$-algèbres
\add{commutatives unitaires}
\[
h : \mathcal{D} \longrightarrow C
\]
rendant commutatif le diagramme

\page{6}

\begin{tikzcd}
B \otimes C & C \arrow[l, "q_C"'] \\
C \arrow[u, "\mathrm{pr}_2"] & \mathcal{D} \arrow[l, "h"] \arrow[u, "h"']
\end{tikzcd}

\struck{i.e. \ill{}}

est-il vrai qu'on peut factoriser $h$ \uncertain{à} l'aide de
$h' : \mathcal{D} \to A$ (une augmentation) par la formule
\[
h(x) = h'(x) . 1_C \qquad ??
\]

La commutativité du diagramme ci-dessus signifie simplement que $h$ prend
ses valeurs dans la sous-algèbre $C'$ de $C$ \struck{formée des $x \in$}
définie par
\[
x \in C' \iff q_C x = 1_B \otimes x
\]
\note{avant $q_C x$, un groupe biffé illisible}
i.e. les \emph{éléments} \add{invariants} (\emph{primitifs}) de $C$ (rel.
à $B$).\note{« invariants », souligné, est ajouté au-dessus de
« primitifs », entouré} $C'$ est une sous-$A$-algèbre de $C$, \ill{} on
voit \uncertain{qu'on a tout ce qu'il faut} (\ill{} si $C$ \ill{} un fibré
principal).\note{lecture de cette ligne très incertaine}

\emph{Le quotient de $C$ par le groupe d'opérateurs $B$ dans la catégorie
$\mathcal{C}$ existe, et est donné par la sous-algèbre $C'$ de $C$ formée
des éléments} \add{invariants} \struck{\ill{}} \emph{primitifs} de $C$
(\uncertain{par.} à $B$).

\page{7}

Mais attention, ce résultat ne dit pas \uncertain{grand-chose} : si on se
place dans la catégorie des \emph{Schémas} sur $\mathrm{Spec}(A)$, il
\ill{} par l'\uncertain{existence} des quotients dans cette catégorie,
\struck{\ill{}} si le quotient y existe, il peut être distinct de celui
\ill{} de la définition. Supposons par ex. $A$ = corps alg. clos,
$C$ = algèbre \ill{} \ill{} \uncertain{linéaire} $E$, $B$ = algèbre
\ill{} d'un sous-groupe de Borel $G$. \uncertain{Alors} $E/G$ dans la
catégorie des schémas existe bien et est une variété complète. Il en
résulte que le quotient dans la catégorie des \uncertain{variétés affines}
est \uncertain{réduit au point}.\note{un grand crochet en marge gauche
embrasse ce paragraphe}

Le cas \uncertain{qui nous intéresse}, \struck{\ill{}}

\marginal{$\|$}
[\emph{$C$ étant principal homogène}] \struck{les éléments} \emph{$C'$ est
isomorphe à $A$} [\emph{par $\lambda \mapsto \lambda 1_C$}].\note{les deux
lignes sont soulignées d'un trait ondulé et marquées d'un double trait
vertical en marge}

En effet, si $x$ est semi-primitif, alors le diagramme commutatif (*)
donne
\[
x \otimes 1_C = 1_C \otimes x \quad\text{dans } C \otimes C
\]

\page{8}

\emph{Supposons que} $\lambda \mapsto \lambda 1_C$ soit un isomorphisme de
$A$ sur un facteur direct du $A$-module $C$ [pour ceci, il suffit, comme on
voit facilement, que $C$ soit un $A$-module \struck{libre de type}
\ill{} \uncertain{ayant une base contenant} $1_C$, \uncertain{par ex.} que
$C$ soit un $A$-module libre de type fini]. Alors la condition précédente
implique facilement $x = \lambda 1_C$, où $\lambda \in A$ est déterminé de
façon unique. Dans ce cas, on a donc bien ce qu'on voulait.

\emph{N.B.} En général, \struck{\ill{}} il ne sera pas vrai que le quotient
de $C$ par $B$ (dans $\mathcal{C}$) soit $A$, si les conditions \ill{}
\uncertain{pas remplies} : par exemple \uncertain{prenant} $C = B/JB$, où
$J$ est un idéal de $A$, \struck{\ill{}} \struck{et définissant} en prenant
la $q_C$ structure d'algèbre quotient, i.e. en définissant
$q_C : C \to B \otimes C$ i.e.
\[
q_C : B/JB \longrightarrow (B \otimes B)/J.(B \otimes B)
\]
en réduisant $p_B$ mod $J$. Les axiomes restent vérifiés, cependant
l'algèbre des primitifs de $C$ est $A/J$, et non $A$\ldots

\page{9}

Supposons maintenant que $C$ soit lui aussi un $A$-module libre (cela
montre que la condition relative aux \uncertain{quasi}-primitifs est une
conséquence de l'autre condition). Je dis que pour que \struck{$\varphi$}
\[
\text{(**)} \qquad \varphi : C \otimes_A C \longrightarrow \mathrm{Hom}_A(\mathcal{U}, C)
\]
\emph{soit un isomorphisme}, \emph{il faut et il suffit que l'hom
analogue défini en $\varphi$}
\[
\psi : C \otimes \mathcal{U} \longrightarrow \mathrm{Hom}(C, C)
\]
\emph{par la condition}
\[
\boxed{\ \psi(x \otimes u) . y = x\, u(y)\ }
\]
\emph{soit un isomorphisme}.\marginal{?}\note{la première lettre de
$C \otimes_A C$ est surchargée ; « (**) » est écrit à gauche de la ligne
de $\psi$, sans qu'on voie auquel des deux homomorphismes il se rapporte}
[Donc, si $B$ et $C$ sont libres de type fini, on a deux caractérisations
\add{équivalentes} \struck{\ill{} \ill{} pour que $C$ soit} du fait que
$C$ est un fibré principal homogène sous \uncertain{$C$}]\note{sic ; on
attend $B$. « équivalentes » est écrit au-dessus du passage biffé}

\struck{Supposons} \struck{en} \struck{effet} \struck{que} \struck{$\varphi$} \struck{soit} \struck{un} \struck{isomorphisme,} \struck{i.e.} \struck{\ill{}} \struck{\ill{}} \struck{:} \struck{$A' = C$,} \struck{donc} \struck{$C^{A'} \otimes_{A'} C^{A'}$} \struck{${}\to{}$} \struck{$\mathrm{Hom}_{A'}(\mathcal{U}^{A'}, C^{A'})$} \struck{\ill{}} \struck{$B^{A'}$} \struck{est} \struck{«} \struck{fibré} \struck{hom.} \struck{principal} \struck{sous} \struck{$B = \mathcal{U}'$} \struck{»} \struck{dans} \struck{la} \struck{catégorie} \struck{des} \struck{$A$-algèbres} \struck{\ill{}.} \struck{Faisant} \struck{le} \struck{produit} \struck{\ill{}} \struck{dans} \struck{la} \struck{catégorie} \struck{des} \struck{$A'$-algèbres} \struck{\ill{}} \struck{le} \struck{fibré} \struck{\ill{}} \struck{$A' = C$,} \struck{\ill{}} \struck{$C^{A'}$} \struck{${}= C \otimes_A A'$} \struck{sur} \struck{$B^{A'}$} \struck{${}= \mathcal{U}' \otimes_A A'$} \struck{${}= (\mathcal{U}^{A'})'$} \struck{est} \struck{isomorphe} \struck{:} \struck{[\uncertain{car}} \struck{on} \struck{a} \struck{pris} \struck{$A' = C$].}\note{toute la moitié inférieure de la page est barrée de
longues diagonales ; lecture très incomplète}

\page{10}

Supposons \struck{à} nouveau que $B$ soit libre de type fini sur $A$,
\struck{alors} et exprimons le fait que $C$ est un fibré principal
\emph{homogène} à l'aide de $\mathcal{U}$, \ill{}. Cela doit donc
exprimer d'abord que l'hom d'algèbres
\[
\varphi : C \otimes C \longrightarrow \mathrm{Hom}_A(\mathcal{U}, C)
\]
[où la structure d'algèbre du deuxième membre est définie par

\begin{tikzcd}
\mathcal{U} \arrow[r, "\Delta_{\mathcal{U}}"] \arrow[rrr, bend right=20, "v * w"'] & \mathcal{U} \otimes \mathcal{U} \arrow[r, "v \otimes w"] & C \otimes C \arrow[r, "\Delta_C"] & C
\end{tikzcd}

]
qui est défini par
\struck{$\varphi(1_C \otimes x) . (u) = \eta(u) x$}
\struck{$\varphi(x \otimes 1_C)(u) = u . x$}
\[
\boxed{\ \varphi(x \otimes y) . u = u(x) . y\ }
\]
est un \emph{isomorphisme} de $C \otimes C$ sur
$\mathrm{Hom}_A(\mathcal{U}, C)$.

De plus, on veut que l'hom $\lambda \mapsto \lambda . 1_C$ soit un
\emph{isomorphisme} de $A$ sur l'ens. des éléments semi-primitifs de $C$,
i.e. tels que
\[
u . x = \eta(u) . x
\]
\note{les deux derniers alinéas sont embrassés par un grand crochet en
marge gauche ; la formule est soulignée d'un trait ondulé}

\page{11}

Catégorie $\mathcal{C}$ avec unité : \ill{} $e$, et produits.

$G$ groupe dans $\mathcal{C}$ \quad $p_G : G \times G \to G$

$E$ $\mathcal{C}$-espace : $\mathcal{C}$-groupe $G$ opérant :
\quad $q_E : G \times E \to E$

On dit que $E$ est formellement \emph{hom. principal} si l'hom
\[
G \times E \xrightarrow{\ (q_E, p_2)\ } E \times E
\]
est un \emph{isomorphisme}.\marginal{$x = g(x, y)\, y$} Cela signifie
aussi que l'on peut trouver un hom
\[
E \times E \xrightarrow{\ g_E\ } G ,
\]
tel que

(i)

\begin{tikzcd}
E \times E \arrow[r, "{(g, p_2)}"] \arrow[rr, bend right=30, "p_1"'] & G \times E \arrow[r, "q_E"] & E
\end{tikzcd}

est commutatif

(ii)

\begin{tikzcd}
G \times E \arrow[r, "{(q_E, p_2)}"] \arrow[rr, bend right=30, "p_1"'] & E \times E \arrow[r, "g_E"] & G
\end{tikzcd}

est commutatif

Si $\mathcal{C}$ est la catégorie des $A$-algèbres (commutatives avec
unité), alors $G$, $E$ correspondent à des $A$-algèbres \struck{\ill{}}
$B$, $C$, et on a
\[
\text{(*)} \qquad
\left\lbrace
\begin{array}{l}
p_B : B \longrightarrow B \otimes B \\
q_C : C \longrightarrow B \otimes C
\end{array}
\right.
\]
et le fait que $E$ soit formellement \emph{un} fibré hom. principal
signifie que
\[
\varphi : C \otimes C \xrightarrow{\ (q_C, p_2)\ } B \otimes C
\]
est un isomorphisme, i.e. qu'il existe
\[
g_C : B \longrightarrow C \otimes C
\]

\page{12}

rendant commutatifs les diagrammes

(i)

\begin{tikzcd}
C \arrow[r, "q_C"] \arrow[rr, bend right=30, "p_1"'] & B \otimes C \arrow[r, "{(g_C, p_2)}"] & C \otimes C
\end{tikzcd}

(ii)

\begin{tikzcd}[column sep=large]
B \arrow[r, "g_C"] \arrow[rr, bend right=30, "p_1"'] & C \otimes C \arrow[r, "{\varphi = (q_C, p_2)}"] & B \otimes C
\end{tikzcd}

D'autre part, \uncertain{on démontre} en définissant
\[
\psi : \mathrm{Hom}_A(B, C) \longrightarrow \mathrm{Hom}_A(C, C)
\]
en posant $\psi(v) = \Delta_C (v \otimes \mathrm{id}_C) q_C$ :

\begin{tikzcd}
C \arrow[r, "q_C"] \arrow[rrr, bend right=20, "\psi(v)"'] & B \otimes C \arrow[r, "v \otimes \mathrm{id}_C"] & C \otimes C \arrow[r, "\Delta_C"] & C
\end{tikzcd}

Je dis que \emph{si} $\varphi$ \emph{est un isomorphisme}, $\psi$
\emph{en est un}. On va en effet construire, utilisant $g_C$, un hom
\[
\psi' : \mathrm{Hom}(C, C) \longrightarrow \mathrm{Hom}(B, C)
\]
qui sera inverse de $\psi$. On pose
$\psi'(w) = \Delta_C (w \otimes \mathrm{id}_C) g_C$ :

\begin{tikzcd}
B \arrow[r, "g_C"] \arrow[rrr, bend right=20, "\psi'(w)"'] & C \otimes C \arrow[r, "w \otimes \mathrm{id}_C"] & C \otimes C \arrow[r, "\Delta_C"] & C
\end{tikzcd}

On doit prouver $\psi'(\psi(v)) = v$, $\psi \psi'(w) = w$. C'est facile :
\struck{$\psi\psi'(w) = {}$} \struck{$\Delta_C\bigl[(\Delta_C (w \otimes \mathrm{id}) g_C)$} \struck{$\otimes \mathrm{id}_C\bigr] q_C$}

\page{13}

1°) \emph{$\psi \psi'(w) = w$}\note{devant « 1° », un signe griffonné}

\begin{tikzcd}[column sep=small, row sep=small, nodes={font=\scriptsize}]
C \arrow[d, "q_C"] & & C \arrow[d, "p_1"] \\
B \otimes C \arrow[d, "g_C \otimes \mathrm{id}_C"] & & C \otimes C \arrow[d, "w \otimes \mathrm{id}_C"] \\
C \otimes C \otimes C \arrow[d, "w \otimes \mathrm{id}_{C \otimes C}"] & & C \otimes C \arrow[d, "\Delta_C"] \\
C \otimes C \otimes C \arrow[d, "\Delta_C \otimes \mathrm{id}_C"] & & C \\
C \otimes C \arrow[d, "\Delta_C"] & & \\
C & &
\end{tikzcd}

\note{colonne de gauche : la première, où les deux flèches de chaque étage
sont écrites côte à côte ($g_C$ et $\mathrm{id}_C$, $w$ et
$\mathrm{id}_{C \otimes C}$, $\Delta_C$ et $\mathrm{id}_C$), les facteurs
concernés marqués d'accolades ; colonne de droite : la dernière des six.
Entre elles, quatre colonnes intermédiaires, reliées par des accolades, où
l'on échange l'ordre des étapes : $C \otimes C \otimes C$ par
$\mathrm{id}_C$ et $\Delta_C$ vers $C \otimes C$ puis $\Delta_C$ ;
$C \otimes C \otimes C$ par $w$ et $\Delta_C$ ; $C \otimes C \otimes C$ par
$\mathrm{id}_C$, $\Delta_C$, puis $w$, $\mathrm{id}_C$ ; $C \to B \otimes C
\xrightarrow{(g_C, p_2)} C \otimes C \to C$. Deux fins de colonne sont
biffées de hachures}

Sur la dernière colonne, on reconnaît aussitôt qu'on a un composé
identique à $w$.

2°) $\psi' \psi(v) = v$

\begin{tikzcd}[column sep=small, row sep=small, nodes={font=\scriptsize}]
B \arrow[d, "g_C"] & & B \arrow[d, "p_1"] \\
C \otimes C \arrow[d, "q_C \otimes \mathrm{id}_C"] & & B \otimes C \arrow[d, "v \otimes \mathrm{id}_C"] \\
B \otimes C \otimes C \arrow[d, "v \otimes \mathrm{id}_{C \otimes C}"] & & C \otimes C \arrow[d, "\Delta_C"] \\
C \otimes C \otimes C \arrow[d, "\Delta_C \otimes \mathrm{id}_C"] & & C \\
C \otimes C \arrow[d, "\Delta_C"] & & \\
C & &
\end{tikzcd}

\note{même disposition : la première et la dernière des six colonnes ; les
intermédiaires passent par $C \otimes C \otimes C \to C \otimes C \to C$,
$B \otimes C \otimes C \xrightarrow{v,\ \Delta_C} B \otimes C$ (sic),
$B \otimes C \otimes C \xrightarrow{\mathrm{id}_B,\ \Delta_C} B \otimes C
\xrightarrow{v,\ \mathrm{id}_C} B \otimes C$ (sic), et
$C \otimes C \xrightarrow{(q_C, p_2)} B \otimes C$ ; les « $B$ » de ces
colonnes intermédiaires, là où l'on attend $C$, sont tels sur la page}

Sur la dernière colonne, on voit aussitôt que le composé est identique à
$v$.

\page{14}

Réciproquement, supposons $\psi$ \emph{un isomorphisme}, \emph{montrons
qu'alors} $\varphi$ \emph{est un isomorphisme}, du moins si on admet que
$C$ admet une base finie sur $A$. Soit en effet $\psi'$ l'isomorphisme
inverse de $\psi$
\[
\psi' : \mathrm{Hom}_A(C, C) \longrightarrow \mathrm{Hom}(B, C)
\]
il s'identifie à un hom
\struck{\ill{}}
\[
\mathrm{Hom}(C, C) \otimes B \longrightarrow C
\]
i.e. à un hom
\[
\psi'_1 : B \longrightarrow \mathrm{Hom}(\mathrm{Hom}(C, C), C)
\]
On veut en tirer un hom $B \to C \otimes C$ en utilisant un hom naturel
\[
\rho : \mathrm{Hom}(\mathrm{Hom}(C, C), C) \longrightarrow C \otimes C
\]
Pour construire ce dernier, on note que le premier est isomorphe :
\[
(C' \otimes C)' \otimes C = C \otimes C' \otimes C
\]
\note{à la suite, un groupe biffé : \struck{\ill{} $\mathrm{Hom}(C, C)$}}
\struck{[En} \struck{général,} \struck{on} \struck{a} \struck{\ill{}} \struck{$P \otimes \mathrm{Hom}(Q, R)$} \struck{${}\to{}$} \struck{$\mathrm{Hom}(\mathrm{Hom}(P, Q), R)$} \struck{i.e.} \struck{$P \otimes \mathrm{Hom}(Q, R)$} \struck{${}\otimes \mathrm{Hom}(P, Q)$} \struck{${}\to R$,} \struck{défini} \struck{par} \struck{$(p, v, u)$} \struck{${}\mapsto (v \circ u)(p)$]},\note{ce crochet est
barré d'un trait en zigzag} or l'élément unité de
\uncertain{$C$} définit un $A$-hom naturel
\[
\varepsilon : C' \longrightarrow A
\]
d'où un hom $\rho = \mathrm{id}_C \otimes \varepsilon \otimes \mathrm{id}_C$.

\page{15}

Il faut prouver alors que, posant
\[
g_C = \rho \circ \psi'_1 : B \longrightarrow C \otimes C
\]
l'hom $g_C$ satisfait les relations (i) et (ii).\marginal{???}

Pour ceci, ----\note{le reste de la page est blanc}

\page{16}

Considérons l'hyperalgèbre
\[
\mathcal{U} = A[X] / X^p A[X]
\]
sur l'anneau $A$ \add{de car. $p$} \add{avec} donnée par
\[
\boxed{\ \Delta X = X \otimes 1 + 1 \otimes X\ }
\]
[cela définit d'abord $A[X] \to A[X] \otimes A[X]$, mais
$\Delta X^p = X^p \otimes 1 + 1 \otimes X^p$ est \uncertain{multiple} dans
l'idéal $A[X] \otimes X^p A[X] + X^p A[X] \otimes A[X]$], l'augmentation
étant définie par \struck{\ill{}} $\boxed{\varepsilon(X) = 0}$, et
l'\,$\vee$ par $\boxed{\check{X} = -X}$.\note{« de car. $p$ » et « avec »
sont ajoutés au-dessus de la ligne, le premier avec un renvoi}

Cette hyperalgèbre est \uncertain{isomorphe} à sa duale, \struck{à
l'aide} on la regarde comme \struck{\ill{} des opérateurs} duale d'un
groupe dans la catégorie des $A$-algèbres commutatives.

Soit $C$ une $A$-algèbre commutative. Se faire opérer l'hyperalgèbre
$\mathcal{U}$ sur $C$ revient à définir l'image $D$ de $X$, i.e. à se
donner un opérateur $D$ $A$-linéaire de $C$. De plus
\begin{itemize}
  \item[(i)] $D^p = 0$ (exprime que $D$ définit une représentation
  d'$A$-algèbre de $A[X]/X^p A[X]$ dans $L_A(C, C)$),
  \item[(ii)] $D(xy) = x Dy + y Dx$ (exprime que cette représentation est
  compatible avec $\Delta$)
\end{itemize}
donc \ill{} \uncertain{trouvé} : la donnée d'une \emph{$A$-dérivation de
puissance $p$-ième nulle}.

On veut exprimer maintenant que $C$ est un $B$-espace
\emph{principal homogène} dans la catégorie des $A$-algèbres commutatives.
On suppose d'abord \struck{que \ill{}} que $C$ \struck{\ill{}} \add{est
libre} un $A$-module libre (ce qui semble d'ailleurs une\note{la lettre
$B$ est surchargée et porte un accent ou un prime}

\page{17}

conséquence de notre hypothèse), et \struck{\ill{}} rang égal nécessairement
$p$. Ceci dit il faut exprimer que les
\[
D^i \quad (0 \leq i \leq p-1) \quad \text{forment une \emph{base} du $C$-module } \mathrm{Hom}_A(C, C)
\]
[où on pose $(\lambda u)(y) = \lambda(u(y))$].

\struck{Soit \ill{}} Supposons par ex. que $C$ admette une $p$-base
\add{sur $A$} (réduite à un élément $u$), donc que
\[
C \simeq A[u]/(u^p - a), \qquad \text{où } a \in A .
\]
\struck{Posons} \struck{$D u$} Alors une $A$-dérivation de $C$ est connue
quand on connaît sa valeur sur $u$, qui peut être choisie arbitrairement.
Définissons alors la $A$-dérivation $D$ dans $C$ par
\[
Du = 1
\]
d'où
\[
Du^n = n u^{n-1} .
\]
On a donc $D^p = 0$, donc $C$ est un $\mathcal{U}$-module. De plus,
\struck{la matrice des} identifions, grâce à la base
$(u^i)_{0 \leq i \leq p-1}$, $\mathrm{Hom}(C, C)$ à $C^p$, par
$v \mapsto (v(u^i))_{0 \leq i \leq p-1}$ [c'est un isom. de $C$-modules].
\struck{Alors} La matrice des opérations des $D^i$
($0 \leq i \leq p-1$) est
\[
\begin{pmatrix}
1 & u & u^2 & u^3 & \cdots & u^{p-1} \\
0 & 1 & 2u & 3u^2 & \cdots & (p-1) u^{p-2} \\
0 & 0 & 2.1 & 3.2.u & \cdots & (p-1)(p-2) u^{p-3} \\
0 & 0 & 0 & 3.2.1 & \cdots & (p-1)(p-2)(p-3) u^{p-4} \\
\vdots & & & & \ddots & \vdots \\
0 & 0 & 0 & 0 & \cdots & (p-1) \cdots 1
\end{pmatrix}
\]
\note{la dernière colonne est séparée des autres par un double trait
vertical ; la ligne des points de suspension est notre}

\page{18}

Son déterminant est égal à
\[
0!\, 1!\, 2! \cdots (p-1)!
\]
qui provient donc d'un élément non nul de $\mathbb{Z}/p\mathbb{Z}$, donc
est inversible dans $C$, i.e. les $D^i$ ($0 \leq i \leq p-1$) forment bien
une base de $\mathrm{Hom}_A(C, C)$ sur $C$. \struck{Donc}

\emph{Proposition}. Soit $C$ une $A$-algèbre ayant une $p$-base réduite à
un élément. Alors $C$ est un espace homogène fibré principal sous
$\mathcal{U} = A[X]/X^p A[X]$ (de la façon explicitée plus
haut).\note{un trait vertical en marge gauche embrasse l'énoncé ; sous
lui, un double trait ondulé}

\emph{Questions}. 1) Soit $C$ comme plus haut. Déterminer les façons
possibles de \struck{\ill{} \ill{}} \add{$C$} \uncertain{définir} la
$C$-hyperalgèbre \struck{Hom} formée $\mathrm{Hom}_A(C, C)$ à partir d'une
sous-$A$-hyperalgèbre $\mathcal{U}$ (libre de rang $p$ sur $A$).
Déterminer les $\mathcal{U}$ \struck{pour lesquels} qui sont isomorphes à
l'algèbre \uncertain{enveloppante} \add{\ill{} $A$, \ill{} \ill{}}
\ill{} d'une $p$-algèbre de Lie \add{\uncertain{telle}} (libre de rang 1,
d'une) $p$-algèbre pour \uncertain{une} puissance $p$-ième
nulle.\note{les ajouts interlinéaires de ces deux dernières lignes sont
serrés et en partie illisibles ; l'ordre de lecture est incertain}

\page{19}

2) Déterminons les classes \uncertain{à} isomorphismes près d'espaces
\struck{\ill{}} principaux homogènes sous $B$.\note{sous la ligne, un
double trait de séparation}

Soit $C$ une $A$-algèbre \add{libre de rang $p$}, munie d'une
$A$-dérivation $D$ telle que $D^p = 0$, et que les $D^i$
($0 \leq i \leq p-1$) forment une \emph{base} de
$\mathrm{Hom}_A(C, C)$ sur $C$. Alors \emph{l'ensemble des éléments de
$C$ annulés par $D$ est $A$}. En effet, tout sous-$A$-module $M$ de $C$
qui est \add{libre de rang $m$ et} \emph{facteur direct} dans $C$,
s'identifie à l'ensemble des zéros de \add{$(M^\circ)^\circ$} l'ens.
$M^\circ$ des $v \in \mathrm{Hom}(C, C)$ qui s'annulent sur $M$
(évident), et $M^\circ$ est un sous-$C$-module libre de rang $p - m$ dans
$\mathrm{Hom}(C, C)$, \struck{\ill{}} qui est facteur direct \ill{} des
sous-$C$-modules ; \ill{} \ill{} \emph{si} \ill{} \ill{} $u_i$ dans
$M^\circ$ \add{forment la base d'un sous-module $L$ de
$\mathrm{Hom}(C, C)$} \struck{libre de \ill{} \ill{}} qui \ill{} facteurs
directs, alors nécessairement $L = M^\circ$, d'où $M = L^\circ$, i.e.
$M$ est l'ens. des éléments de $C$ annulés par les $u_i$. \ill{}
\uncertain{applique} ceci\note{en marge gauche, deux notes tournées,
reliées au texte par une accolade : « \ill{} \ill{} sous-$A$-modules
\ill{} $A$ : $C$ » et, plus bas, « Lemme \ill{} » ; lecture très
incertaine}

\page{20}

: $M = A$, et on a $u_i = D^i$ ($1 \leq i \leq p-1$) \ill{}, etc ----

\struck{S'il existe des} \struck{Corollaire} \emph{Il existe dans $C$ un
élément} \add{$u$} \struck{\ill{}} \emph{tel que $Du = 1$, et $u$ est
unique à addition près d'un élément de $A$.}\note{un trait vertical en
marge gauche embrasse l'énoncé}

L'unicité résulte du fait que \struck{\ill{} \ill{} \ill{}}
$D(u - u') = 0$ \struck{\ill{}} \uncertain{implique} $u - u' \in A$.

Montrons enfin que l'on peut trouver $u$ avec $Du = 1$.
\struck{En effet, la question étant locale, on} \ill{}
$A$ local, l'image de $D$ est donc un sous-$A$-module de $C$,
\uncertain{isomorphe} à $C / \operatorname{Ker} D \simeq C/A$, donc
\emph{libre} de rang $p - 1$. Je dis qu'il contient $A$, i.e. que
$D(C)^\circ$ \uncertain{contenu} dans $A^\circ$, i.e. que si
$v = \sum_{0 \leq i \leq p-1} \lambda_i D^i$ est nul sur $D(C)$, alors
$\lambda_0 = 0$. Or l'hypothèse sur $v$ signifie que $vD = 0$, i.e.
$\lambda_0 = \lambda_1 = \cdots = \lambda_{p-2} = 0$, OK.\note{tout ce
passage, depuis « $A$ local », est entouré d'un cadre et traversé de deux
longues obliques ; il n'est pas clair qu'il soit annulé}

\uncertain{Si} $x \in C$, on a $x^p \in A$ \uncertain{si} $x \in C$ car
$Dx^p = 0$.

\emph{Si} $u$ est tel que $Du = 1$, alors les $u^i$
($0 \leq i \leq p-1$) \emph{forment une base} de $C$ sur $A$.
\struck{\ill{} \ill{} \ill{} est locale, \ill{}}

\marginal{En effet, il existe \add{tel} un $A$-endomorphisme $\Theta$ de
$C$ tel que $\Theta(1) = 0$ et $\Theta(u) = 1$ (\ill{} $C$ \ill{}). On a
$\Theta = \sum \lambda_i D^i = \ill{}$ \ill{} $Du = 1$, \ill{}
$u = \Theta \ill{}$, $1 \leq i \leq p-1$}\note{note écrite dans la marge
de gauche, tournée d'un quart de tour ; elle se rapporte à la preuve de
l'existence de $u$ ; lecture très incomplète. La suite, « $A$ est local.
Montrons d'abord que les $u^i$ \ldots\ sont libres sur $A$ », ouvre la
page 21 (lot 2)}

\end{document}
